%% kapa: exam=1 round=23 year=2012 period=1 form=A subject=economics subject_ko=경제원론 seq=10 number=50
%% source: Q-Net 감정평가사 기출문제 (cst00309) artlSeq=5052578 file=raw/1008/5052578/2012년 제23회 감정평가사 제1차시험 기출문제(1교시 A형).pdf
%% conversion: kordoc 4.24.51 -> tools/gen1.py (auto); review=proofread-sonnet; flags=none
\begin{question}{50}
단위당 동일한 종량세율로 생산자 또는 소비자에게 부과하는 조세에 관한 설명으로 옳지 않은 것은?
\begin{choices}
  \item 생산자에게 부과할 때와 소비자에게 부과할 때의 경제적 순손실(deadweight loss)은 같다.
  \item 조세부담의 귀착(tax incidence)은 조세당국과 생산자 및 소비자 간의 협상능력에 의존한다.
  \item 수요의 가격탄력성이 클수록 생산자의 조세부담이 커진다.
  \item 수요의 가격탄력성이 공급의 가격탄력성보다 클수록 생산자의 조세부담분이 커진다.
  \item 수요의 가격탄력성이 0인 재화에 조세를 부과해도 사회후생은 감소하지 않는다.
\end{choices}
\end{question}
\answer{2}
